\documentclass[aspectratio=169]{beamer}

% Theme selection
\usetheme{Madrid}
\usecolortheme{whale}

% --- UNIVERSAL PREAMBLE BLOCK ---
\usepackage{fontspec}
\usepackage[english, bidi=basic, provide=*]{babel}

\babelprovide[import, onchar=ids fonts]{english}

% Set default/Latin font to Sans Serif (Noto Sans)
\babelfont{rm}{Noto Sans}
\babelfont{sf}{Noto Sans} % Beamer primarily uses the sans-serif family
\babelfont{tt}{Noto Sans Mono}

% Remove navigation symbols for a cleaner look
\setbeamertemplate{navigation symbols}{}

% Title slide information
\title[NSO Nepal AI Guide]{Using AI Tools to Prepare \\ Official Statistical Reports}
\subtitle{A Practical Learning Guide for NSO Nepal Staff}
\author[NSO Nepal]{National Statistics Office of Nepal}
\date{June 2026}

\begin{document}

% Slide 1: Title Slide
\begin{frame}
    \titlepage
\end{frame}

% Slide 2: What AI Is (and Isn't)
\begin{frame}{Step 1: The Right Mental Model}
    \begin{block}{What is AI?}
        Think of AI tools (ChatGPT, Claude, Gemini, DeepSeek, Qwen) like a very well-read, fast-writing colleague.
    \end{block}
    
    \vspace{0.5cm}
    
    \textbf{What AI is EXCELLENT for:}
    \begin{itemize}
        \item Drafting narrative paragraphs from data you provide.
        \item Rewriting awkward sentences into formal English.
        \item Generating report outlines and structures.
    \end{itemize}
    
    \vspace{0.5cm}
    
    \textbf{What AI is UNRELIABLE for:}
    \begin{itemize}
        \item Quoting specific statistics not provided to it.
        \item AI can "hallucinate" (invent) plausible-sounding but incorrect numbers.
    \end{itemize}
\end{frame}

% Slide 3: Prompt Engineering
\begin{frame}{Step 2: Prompt Engineering (The RCTF Framework)}
    \textbf{The quality of AI output depends entirely on your prompt.}
    
    \vspace{0.5cm}
    
    Use the \textbf{RCTF Framework} for NSO reports:
    \begin{table}[]
        \centering
        \renewcommand{\arraystretch}{1.5}
        \begin{tabular}{|l|p{8cm}|}
            \hline
            \textbf{R}ole & Tell AI who it is (\textit{"You are an official report writer for NSO Nepal"}) \\ \hline
            \textbf{C}ontext & Provide your data (\textit{"Here is our CPI data for Q2 2024..."}) \\ \hline
            \textbf{T}ask & State exactly what you want (\textit{"Write a press release narrative"}) \\ \hline
            \textbf{F}ormat & Specify length, tone (\textit{"Formal, 150 words, no bullet points"}) \\ \hline
        \end{tabular}
    \end{table}
\end{frame}

% Slide 3b: RCTF Example
\begin{frame}{Step 2: RCTF Framework in Action (Example)}
    \begin{exampleblock}{Full Prompt Example for a CPI Bulletin}
        \textbf{[R]ole:} You are a senior economist drafting text for NSO Nepal's Consumer Price Index (CPI) bulletin. \\
        
        \vspace{0.2cm}
        
        \textbf{[C]ontext:} Data (FY 2023/24, Nepal): \\
        - Overall CPI: 5.4\% (year-on-year) \\
        - Food and beverages: 6.1\% \\
        - Non-food and services: 4.9\% \\
        
        \vspace{0.2cm}
        
        \textbf{[T]ask:} Write the opening paragraph for the monthly CPI press release. State the headline figure and compare food vs. non-food inflation. \\
        
        \vspace{0.2cm}
        
        \textbf{[F]ormat:} Use a formal, official tone appropriate for a government press release. Be no longer than 120 words. Do NOT fabricate any additional statistics.
    \end{exampleblock}
\end{frame}

% Slide 4: Outlining and Narratives
\begin{frame}{Steps 3 \& 4: Outlines and Data-to-Narrative}
    \begin{columns}[T]
        \begin{column}{0.48\textwidth}
            \textbf{Structuring Reports}
            \begin{itemize}
                \item Never start with a blank page.
                \item Ask AI to generate a complete section-by-section outline.
                \item Adjust the outline to match NSO house style.
            \end{itemize}
        \end{column}
        
        \begin{column}{0.48\textwidth}
            \textbf{Drafting Narratives}
            \begin{enumerate}
                \item \textbf{Extract data} (Excel/SPSS).
                \item \textbf{Paste data IN the prompt}.
                \item \textbf{Specify highlights} (tell AI what trends matter).
                \item \textbf{Request narrative} (RCTF).
                \item \textbf{Verify every number}.
            \end{enumerate}
        \end{column}
    \end{columns}
\end{frame}

% Slide 5: Editing and Tone
\begin{frame}{Steps 5 \& 6: Executive Summaries \& Editing}
    \textbf{Drafting Summaries \& Press Releases}
    \begin{itemize}
        \item AI is excellent at translating complex statistical findings into clear language for policymakers and media.
        \item Always include a clear headline and key findings.
    \end{itemize}
    
    \vspace{0.5cm}
    
    \textbf{AI as an Editor}
    \begin{itemize}
        \item Paste your draft and ask AI to improve clarity, formality, and consistency.
        \item \textit{Example Prompt:} "Rewrite this in formal, official language without losing meaning. Do NOT change any statistics."
    \end{itemize}
\end{frame}

% Slide 6: Fact-Checking
\begin{frame}{Step 7: Fact-Checking and Verification}
    \begin{alertblock}{Never Trust AI Blindly}
        AI can subtly introduce errors by rounding numbers, switching units, or confusing time periods.
    \end{alertblock}
    
    \vspace{0.3cm}
    \textbf{The NSO Verification Checklist:}
    \begin{itemize}
        \item [\checkmark] Every number verified against the source table.
        \item [\checkmark] Units are correct (\%, percentage points, NRs, millions).
        \item [\checkmark] Time periods match (e.g., FY 2023/24).
        \item [\checkmark] Geographic scope is correct (national/urban/rural).
        \item [\checkmark] No invented statistics not present in original data.
    \end{itemize}
\end{frame}

% Slide 7: Ethics and Privacy
\begin{frame}{Step 8: Ethics, Privacy, and Responsible AI}
    \begin{block}{The Golden Question}
        Before pasting any data into an AI tool, ask: \\
        \textbf{"Would I post this on our website?"}
    \end{block}
    
    \vspace{0.5cm}
    \begin{columns}[T]
        \begin{column}{0.48\textwidth}
            \textcolor{green}{\textbf{✓ Safe to Paste}}
            \begin{itemize}
                \item Published aggregate statistics
                \item Your own draft text for editing
                \item Publicly available NSO reports
            \end{itemize}
        \end{column}
        
        \begin{column}{0.48\textwidth}
            \textcolor{red}{\textbf{X NEVER Paste}}
            \begin{itemize}
                \item Individual/household microdata
                \item Embargoed pre-release data
                \item Internal HR or policy docs
            \end{itemize}
        \end{column}
    \end{columns}
\end{frame}

% Slide 8: The Five Golden Rules
\begin{frame}{Summary: The Five Golden Rules}
    \begin{enumerate}
        \Large
        \item \textbf{You bring the data. AI brings the words.}
        \vspace{0.3cm}
        \item \textbf{Always verify every number in AI output.}
        \vspace{0.3cm}
        \item \textbf{Never paste confidential or unpublished data into AI.}
        \vspace{0.3cm}
        \item \textbf{Use the RCTF framework for every prompt.}
        \vspace{0.3cm}
        \item \textbf{AI is a tool. You are the expert.}
    \end{enumerate}
\end{frame}

\end{document}